% Options for packages loaded elsewhere
% Options for packages loaded elsewhere
\PassOptionsToPackage{unicode}{hyperref}
\PassOptionsToPackage{hyphens}{url}
\PassOptionsToPackage{dvipsnames,svgnames,x11names}{xcolor}
%
\documentclass[
  letterpaper,
  DIV=11,
  numbers=noendperiod]{scrartcl}
\usepackage{xcolor}
\usepackage{amsmath,amssymb}
\setcounter{secnumdepth}{-\maxdimen} % remove section numbering
\usepackage{iftex}
\ifPDFTeX
  \usepackage[T1]{fontenc}
  \usepackage[utf8]{inputenc}
  \usepackage{textcomp} % provide euro and other symbols
\else % if luatex or xetex
  \usepackage{unicode-math} % this also loads fontspec
  \defaultfontfeatures{Scale=MatchLowercase}
  \defaultfontfeatures[\rmfamily]{Ligatures=TeX,Scale=1}
\fi
\usepackage{lmodern}
\ifPDFTeX\else
  % xetex/luatex font selection
\fi
% Use upquote if available, for straight quotes in verbatim environments
\IfFileExists{upquote.sty}{\usepackage{upquote}}{}
\IfFileExists{microtype.sty}{% use microtype if available
  \usepackage[]{microtype}
  \UseMicrotypeSet[protrusion]{basicmath} % disable protrusion for tt fonts
}{}
\makeatletter
\@ifundefined{KOMAClassName}{% if non-KOMA class
  \IfFileExists{parskip.sty}{%
    \usepackage{parskip}
  }{% else
    \setlength{\parindent}{0pt}
    \setlength{\parskip}{6pt plus 2pt minus 1pt}}
}{% if KOMA class
  \KOMAoptions{parskip=half}}
\makeatother
% Make \paragraph and \subparagraph free-standing
\makeatletter
\ifx\paragraph\undefined\else
  \let\oldparagraph\paragraph
  \renewcommand{\paragraph}{
    \@ifstar
      \xxxParagraphStar
      \xxxParagraphNoStar
  }
  \newcommand{\xxxParagraphStar}[1]{\oldparagraph*{#1}\mbox{}}
  \newcommand{\xxxParagraphNoStar}[1]{\oldparagraph{#1}\mbox{}}
\fi
\ifx\subparagraph\undefined\else
  \let\oldsubparagraph\subparagraph
  \renewcommand{\subparagraph}{
    \@ifstar
      \xxxSubParagraphStar
      \xxxSubParagraphNoStar
  }
  \newcommand{\xxxSubParagraphStar}[1]{\oldsubparagraph*{#1}\mbox{}}
  \newcommand{\xxxSubParagraphNoStar}[1]{\oldsubparagraph{#1}\mbox{}}
\fi
\makeatother


\usepackage{longtable,booktabs,array}
\usepackage{calc} % for calculating minipage widths
% Correct order of tables after \paragraph or \subparagraph
\usepackage{etoolbox}
\makeatletter
\patchcmd\longtable{\par}{\if@noskipsec\mbox{}\fi\par}{}{}
\makeatother
% Allow footnotes in longtable head/foot
\IfFileExists{footnotehyper.sty}{\usepackage{footnotehyper}}{\usepackage{footnote}}
\makesavenoteenv{longtable}
\usepackage{graphicx}
\makeatletter
\newsavebox\pandoc@box
\newcommand*\pandocbounded[1]{% scales image to fit in text height/width
  \sbox\pandoc@box{#1}%
  \Gscale@div\@tempa{\textheight}{\dimexpr\ht\pandoc@box+\dp\pandoc@box\relax}%
  \Gscale@div\@tempb{\linewidth}{\wd\pandoc@box}%
  \ifdim\@tempb\p@<\@tempa\p@\let\@tempa\@tempb\fi% select the smaller of both
  \ifdim\@tempa\p@<\p@\scalebox{\@tempa}{\usebox\pandoc@box}%
  \else\usebox{\pandoc@box}%
  \fi%
}
% Set default figure placement to htbp
\def\fps@figure{htbp}
\makeatother





\setlength{\emergencystretch}{3em} % prevent overfull lines

\providecommand{\tightlist}{%
  \setlength{\itemsep}{0pt}\setlength{\parskip}{0pt}}



 


\KOMAoption{captions}{tableheading}
\makeatletter
\@ifpackageloaded{caption}{}{\usepackage{caption}}
\AtBeginDocument{%
\ifdefined\contentsname
  \renewcommand*\contentsname{Table of contents}
\else
  \newcommand\contentsname{Table of contents}
\fi
\ifdefined\listfigurename
  \renewcommand*\listfigurename{List of Figures}
\else
  \newcommand\listfigurename{List of Figures}
\fi
\ifdefined\listtablename
  \renewcommand*\listtablename{List of Tables}
\else
  \newcommand\listtablename{List of Tables}
\fi
\ifdefined\figurename
  \renewcommand*\figurename{Figure}
\else
  \newcommand\figurename{Figure}
\fi
\ifdefined\tablename
  \renewcommand*\tablename{Table}
\else
  \newcommand\tablename{Table}
\fi
}
\@ifpackageloaded{float}{}{\usepackage{float}}
\floatstyle{ruled}
\@ifundefined{c@chapter}{\newfloat{codelisting}{h}{lop}}{\newfloat{codelisting}{h}{lop}[chapter]}
\floatname{codelisting}{Listing}
\newcommand*\listoflistings{\listof{codelisting}{List of Listings}}
\makeatother
\makeatletter
\makeatother
\makeatletter
\@ifpackageloaded{caption}{}{\usepackage{caption}}
\@ifpackageloaded{subcaption}{}{\usepackage{subcaption}}
\makeatother
\usepackage{bookmark}
\IfFileExists{xurl.sty}{\usepackage{xurl}}{} % add URL line breaks if available
\urlstyle{same}
\hypersetup{
  pdftitle={Los fundamentos lógica y demostraciones},
  colorlinks=true,
  linkcolor={blue},
  filecolor={Maroon},
  citecolor={Blue},
  urlcolor={Blue},
  pdfcreator={LaTeX via pandoc}}


\title{Los fundamentos lógica y demostraciones}
\author{}
\date{}
\begin{document}
\maketitle


La lógica es la base de todo razonamiento matemático y de todo
razonamiento automatizado. Para comprender las matemáticas debemos
entender en que consiste un razonamiento matemático correcto, es decir,
una desmotración. Una vez que demostremos que una afirmación matemática
es cierta, la llamamos \textbf{teorema}.

Las demostraciones se utilizan para verificar que los programas
informáticos generen el resultado correcto para todos los valores de
entrada posible.

\section{Lógica proposicional}\label{luxf3gica-proposicional}

Las reglas de la lógica dotan de un significado preciso a los enunciados
matemáticos. Estas reglas se utilizan para distinguir entre argumentos
matemáticos válido o inválidos.

\section{Proposiciones}\label{proposiciones}

Una proposición es una oración declarativa (es decir, una oración que
afirma un hecho) que es verdadera o falsa, pero no ambas cosa a la vez.

Todas las oraciones declarativas son proposiciones.

\begin{enumerate}
\def\labelenumi{\arabic{enumi}.}
\tightlist
\item
  Washington, D.C, es la capital de estados unidos.
\item
  Toronto es la capital de Canada.
\item
  1 + 1 = 2.
\item
  2 + 2 = 3.
\end{enumerate}

Considera las siguientes frases.

\begin{enumerate}
\def\labelenumi{\arabic{enumi}.}
\tightlist
\item
  ¿Qué hora es?
\item
  Lee esto con atención.
\item
  x + 1 = 2.
\item
  x + y = z.
\end{enumerate}

La oraciones 1 y 2 no son proposiciones porque no son oraciones
declarativas. Las oraciones 3 y 4 no son proposiciones porque no son ni
verdaderas ni falsas.

Utilizamos letras para designar \textbf{variables proposicionales}, es
decir variables que representan proposiciones. Las letras convencionales
utilizadas para las variables proposicionales son p, q, r, s, el valor
de verdad de un proposición es denotador por \textbf{T}, si la
proposición es verdadera, y el valor de verdad de una proposición es
falso, denotado por una \textbf{F}, si una proposición es falsa.

El área de la lógica que se ocupa de las proposiciones se denomina
lógica proposicional. Muchas afirmaciones matemáticas se construyen
combinando una o más proposiciones.

\subsection{Definición 1.}\label{definiciuxf3n-1.}

Sea \(p\) una proposición. La negación de \(p\), denotadada por
\(\neg p\) es el enunciado.

\textbf{no es cierto que p.}

La proposición \(\neg p\) se lee no p.~El valor de verdad de la negación
p, \(\neg p\), es el opuesto del valor de verdad de p.

\textbf{Conectivos:} Son aquellos operadores lógicos que se utilizan
para formar nuevas proposiciones a partir de dos o más proposiciones
existentes.

\begin{longtable}[]{@{}ll@{}}
\toprule\noalign{}
p & \(\neg p\) \\
\midrule\noalign{}
\endhead
\bottomrule\noalign{}
\endlastfoot
V & F \\
F & V \\
\end{longtable}

\subsection{Definición 2.}\label{definiciuxf3n-2.}

Sean \(p\) y \(q\) proposiciones. La conjunción de \(p\) y \(q\),
denotada por \(p \land q\), que es la proposición \(p\) y \(q\). La
conjunción \(p \land q\) es verdadera cuando tanto p como q son
verdaderas, y falsa en cualquier otro caso.

\begin{longtable}[]{@{}lll@{}}
\toprule\noalign{}
p & q & \(p \land q\) \\
\midrule\noalign{}
\endhead
\bottomrule\noalign{}
\endlastfoot
V & V & V \\
V & F & F \\
F & V & F \\
F & F & F \\
\end{longtable}

\subsection{Definición 3.}\label{definiciuxf3n-3.}

Sean \(p\) y \(q\) proposiciones. La disyunción de \(p\) y \(q\),
denotada por \(p \lor q\), es la proposicón p o q. La disyunción
\(p \lor q\) es falsa cuando tanto p como q son falsas, y es verdadera
en cualquier otro caso.

\begin{longtable}[]{@{}lll@{}}
\toprule\noalign{}
p & q & \(p \lor q\) \\
\midrule\noalign{}
\endhead
\bottomrule\noalign{}
\endlastfoot
V & V & V \\
V & F & V \\
F & V & V \\
F & F & F \\
\end{longtable}

\subsection{Definición 4.}\label{definiciuxf3n-4.}

Sean \(p\) y \(q\) proposiciones. El o exclusivo de \(p\) y \(q\),
denotado por \(p \oplus q\), es la proposición que es verdadera cuando
cuando exactamente una de p y q es verdadera y falsa en cualquier otro
caso.

\begin{longtable}[]{@{}lll@{}}
\toprule\noalign{}
p & q & \(p \oplus q\) \\
\midrule\noalign{}
\endhead
\bottomrule\noalign{}
\endlastfoot
V & V & F \\
V & F & V \\
F & V & V \\
F & F & F \\
\end{longtable}

\section{Enuncidados condicionales}\label{enuncidados-condicionales}

\subsection{Definición 5.}\label{definiciuxf3n-5.}

Sean \(p\) y \(q\) proposiciones. La proposición condicional \(p \to q\)
es la proposición si p, entonces q. La proposición condicional
\(p \to q\) es falsa cuando p es verdadera y q es falsa, y verdadera en
cualquier otro caso. En la proposición condicional \(p \to q\), p se
denomina hipótesis (o antecedente o prmisa) y q se denomia conclusión (o
consecuencia).

La proposición \(p \to q\) se denomina condicional porque que
\(p \to q\) afirma que q verdadera con la condición de que p es cierta.
Una proposicón condicional se denomina implicación.

Una forma util de entender el valor de verdad de la proposición
condicional es pensar en una obligación o un contrato. Por ejemplo, la
promesa que hacen muchos políticos cuando se presentan las elecciones es
``Si salgo elegido, bajaré los impuesto.''

\begin{longtable}[]{@{}lll@{}}
\toprule\noalign{}
p & q & \(p \to q\) \\
\midrule\noalign{}
\endhead
\bottomrule\noalign{}
\endlastfoot
V & V & V \\
V & F & F \\
F & V & V \\
F & F & V \\
\end{longtable}

\subsection{Converso, Contrapositivo e
inverso.}\label{converso-contrapositivo-e-inverso.}

A partir de la proposición condicional \(p \to q\) podemos formar nuevas
proposiciones condicionales. En concreto, hay tres proposiciones
condicionales relacionadas que aperecen con tanta frecuencia que tienen
nombres propios. La proposición \(q \to p\) se denomina
\textbf{reciproca} de \(p \to q\). La \textbf{contrapositiva} de
\(p \to q\) es la proposicón \(\neg q \to \neg p\). La proposicón
\(\neg p \to \neg q\) se le llama inversa de \(p \to q\), solo la
contrapositiva tiene el mismo valor de verdad que \(p \to q\).

Cuando dos proposiciones compuestas tienen siempre el mismo valor de
verdad, las llamamos \textbf{equivalentes}.

\section{Bicondicionales}\label{bicondicionales}

\subsection{Definición 6.}\label{definiciuxf3n-6.}

Sean \(p\) y \(q\) proposiciones. La afirmación bicondicional
\(p \iff q\) es la proposición \(p\) si y solo si \(q\). La proposición
bicondicional \(p \iff q\) es verdadera cuando \(p\) y \(q\) tienen los
mismos valores de verdad, y es falsa en caso contrario. Las afirmaciones
bicondicionales también se denominan biimplicaciones.

Observese que la proposición \(p \iff q\) es verdadera cuando ambas
proposiciones condicionales \(p \to q\) y \(q \to p\) son verdaderas, y
es falsa en cualquier otro caso.

\begin{longtable}[]{@{}lll@{}}
\toprule\noalign{}
p & q & \(p \iff q\) \\
\midrule\noalign{}
\endhead
\bottomrule\noalign{}
\endlastfoot
V & V & V \\
V & F & F \\
F & V & F \\
F & F & V \\
\end{longtable}

\section{Tablas de verdad de proposiciones
compuestas}\label{tablas-de-verdad-de-proposiciones-compuestas}

Utilizando conectores (conjunciones, disyunciones, condicionales,
bicondicionales, \ldots) podemos formarm proposiciones compuestas
complejas que incluyan cualquier número de variables proposicionales.
Podemos utilizar tablas de verdad para determinar los valores de verdad
de estas proposiciones compuestas.

\[
\begin{array}{c|c|c|c|c|c}
p & q & \neg q & p \lor \neg q & p \land q & (p \lor \neg q) \rightarrow (p \land q) \\
\hline
T & T & F & T & T & T \\
T & F & T & T & F & F \\
F & T & F & F & F & T \\
F & F & T & T & F & F \\
\end{array}
\]




\end{document}
